\documentclass[12pt, a4paper]{article}
\usepackage[utf8]{inputenc}
\usepackage[T1]{fontenc}
\usepackage{geometry}
\geometry{left=2.5cm, right=2.5cm, top=2.5cm, bottom=2.5cm}
\usepackage{tikz}
\usetikzlibrary{decorations.pathreplacing}
\usepackage{graphicx}
\usepackage{amsmath, amssymb, amsthm}
\usepackage{setspace}
\usepackage{hyperref}
\usepackage[
    bibstyle=numeric,
    citestyle=authoryear,
    backend=biber,
    natbib=true,
    sorting=nyt]{biblatex}
\usepackage{booktabs}
\usepackage{pgfgantt}
\usepackage{titlesec}
\usepackage{indentfirst}
\addbibresource{reference.bib}

\theoremstyle{definition}
\newtheorem{definition}{Definition}

\title{Model}
\author{}
\date{}

\begin{document}
\maketitle

\section{Model}

\subsection{Environment}

Consider a metropolitan area consisting of a finite set of locations $S = \{1, 2, \ldots, n\}$. These locations are partitioned into municipalities $g \in G$:
\[
S = \bigcup_{g \in G} S^g, \qquad S^g \cap S^{g'} = \emptyset \quad \text{for } g \neq g'.
\]

Each location $i \in S$ is endowed with a fixed supply of land $\bar{H}_i$, allocated between residential and commercial use by local zoning:
\[
\bar{H}_i^R + \bar{H}_i^F = \bar{H}_i,
\]
where $\bar{H}_i^R$ is residential floor space capacity and $\bar{H}_i^F$ is commercial floor space capacity. Both are chosen by the municipal government and assumed to be fully utilized in equilibrium.

Locations are connected by a transportation network. The generalized commuting cost from residence $i$ to workplace $j$ is
\[
\tau_{ij} = \tau_{ij}(I_{ij},\, n_{ij}),
\]
where $I_{ij}$ is transportation investment on link $(i,j)$ and $n_{ij}$ is the equilibrium commuter flow on that link. Investment reduces commuting costs and congestion raises them:
\[
\frac{\partial \tau_{ij}}{\partial I_{ij}} \leq 0, \qquad \frac{\partial \tau_{ij}}{\partial n_{ij}} \geq 0.
\]
Following \cite{Fajgelbaum2020-yv}, I parametrize the commuting cost as
\[
\tau_{ij} = \delta_{ij}^{\tau}\, \frac{n_{ij}^{\rho}}{I_{ij}^{\gamma}},
\]
where $\delta_{ij}^{\tau} > 0$ captures geographic features of link $(i,j)$ and $\rho, \gamma > 0$.

There are three types of agents: households, competitive firms, and governments. Households choose residence and workplace locations. Firms produce a freely tradable good using labor and commercial floor space. Government decisions operate at two levels: municipal governments choose local zoning, and a metropolitan transportation authority chooses investment across all links. In the decentralized equilibrium, municipalities and the transportation authority choose policies simultaneously, each treating the other's decisions as given.

\subsection{Households}

A household chooses a residence location $i$ and a workplace $j$ to maximize utility. Preferences over the numaire good $Q_{ij}$ and residential floor space $h_{ij}^R$ are Cobb-Douglas:
\[
u_{ij}(\omega) = \frac{B_i}{\tau_{ij}}
\left(\frac{Q_{ij}}{\alpha}\right)^{\!\alpha}
\left(\frac{h_{ij}^R}{1-\alpha}\right)^{\!1-\alpha}
v_{ij}(\omega),
\]
where $B_i > 0$ is a residential amenity at location $i$, $\alpha \in (0,1)$ is the expenditure share on the composite good, and $v_{ij}(\omega)$ is an idiosyncratic preference shock drawn independently for each $(i,j)$ pair.

The household's budget constraint is
\[
Q_{ij} + q_i^R h_{ij}^R \leq w_j ,
\]
where $q_i^R$ is the residential rent at location $i$, $w_j$ is the wage at workplace $j$, and $\phi(N_i, L_i)$ is a monetary disutility from residing at $i$. The disutility captures two sources of local negative externalities: residential congestion from neighbors and the presence of workers near the residential location:
\[
\frac{\partial \phi}{\partial N_i} > 0, \qquad \frac{\partial \phi}{\partial L_i} > 0,
\]
where $N_i$ denotes residential population and $L_i$ denotes employment at location $i$ (defined below). I parametrize the disutility as
\begin{align}
    \phi(N_i, L_i) = \eta_N N_i^{1+\rho} + \eta_L L_i^{1+\rho},
\end{align}
where $\eta_N, \eta_L > 0$ are congestion weights on residential crowding and worker influx respectively, and $\rho > 0$ governs the convexity of congestion costs. The marginal congestion costs are $\partial\phi/\partial N_i = (1+\rho)\eta_N N_i^{\rho}$ and $\partial\phi/\partial L_i = (1+\rho)\eta_L L_i^{\rho}$, both increasing in their arguments. Convexity ($\rho > 0$) ensures the municipal first-order condition has an interior solution: as the municipality loosens residential zoning, each additional resident raises congestion by more than the last, so the congestion cost eventually dominates the rent benefit. The ratio $\eta_N/\eta_L$ governs the relative weight incumbents place on residential crowding versus worker influx, and is disciplined by the observed mix of residential and commercial zoning restrictions across locations. The overall scale $(\eta_N, \eta_L)$ is calibrated to match the housing supply elasticity estimates of \citet{Baum-Snow2024-ii}, which proxy for how tightly the municipality's zoning constraint binds at the equilibrium allocation.

Optimal demands are $Q_{ij}^* = \alpha Y_{ij}$ and $h_{ij}^{R*} = (1-\alpha)Y_{ij}/q_i^R$. Substituting into the utility function yields the deterministic component of indirect utility:
\[
\tilde{V}_{ij} = \frac{B_i\, w_{j}}{\tau_{ij}\,(q_i^R)^{1-\alpha}}.
\]

The idiosyncratic shock $v_{ij}(\omega)$ is drawn from a Fréchet distribution with $\Pr(v_{ij} \leq z) = \exp(-T_i E_j z^{-\epsilon})$, where $T_i > 0$ and $E_j > 0$ are location-specific scale parameters and $\epsilon > 1$ governs the dispersion of idiosyncratic tastes. By the Fréchet aggregation property, the probability that a household chooses residence $i$ and workplace $j$ is
\begin{align}
    \pi_{ij} = \frac{T_i E_j\,\tilde{V}_{ij}^{\epsilon}}{\Phi}, \qquad
    \Phi \equiv \sum_{i'}\Phi_{i'}, \qquad
    \Phi_i \equiv \sum_j T_i E_j\,\tilde{V}_{ij}^{\epsilon},
\end{align}
and expected utility is
\begin{align}
    U = \Gamma\!\left(1 - \tfrac{1}{\epsilon}\right)\Phi^{1/\epsilon}.
\end{align}

Let $N$ denote total metropolitan population, fixed in a closed economy. Equilibrium residential population at each location is
\[
N_i = N \sum_j \pi_{ij}.
\]
Residential floor space market clearing at location $i$ pins down the equilibrium rent:
\begin{align}
    q_i^R = \frac{(1-\alpha)\sum_j N\pi_{ij}\,Y_{ij}}{\bar{H}_i^R}.
\end{align}

\subsection{Firms}

At each location $j$, competitive firms produce a freely tradable good using labor $L_j$ and commercial floor space $H_j^F$:
\[
Y_j = A_j L_j^{\gamma}(H_j^F)^{1-\gamma},
\]
where $A_j > 0$ is location-specific productivity and $\gamma \in (0,1)$. Taking wages $w_j$ and commercial rents $q_j^F$ as given, profit maximization yields:
\[
w_j = \gamma A_j \!\left(\frac{H_j^F}{L_j}\right)^{\!1-\gamma}\!, \qquad
q_j^F = (1-\gamma) A_j \!\left(\frac{L_j}{H_j^F}\right)^{\!\gamma}\!.
\]
Commercial floor space is constrained by zoning: $H_j^F = \bar{H}_j^F$. Labor market clearing at $j$ gives
\[
L_j = N \sum_i \pi_{ij},
\]
determined by household location choices. The disutility $\phi(N_i, L_i)$ uses $L_i$ to denote employment at location $i$ acting as a workplace, which is this same quantity evaluated at $j = i$.

\subsection{Municipal Zoning}

Each municipality $g \in G$ chooses the land-use allocation $\{\bar{H}_i^R, \bar{H}_i^F\}_{i \in S^g}$ to maximize the total welfare of its residents. In line with the political economy of local land use \citep{Mast2024-zi, Rollet2025-mk}, the municipal government acts as a representative of its residents, who face a tradeoff between lower housing costs — which favor denser zoning — and lower congestion and fewer workers near their homes — which favor restrictive zoning. Formally, municipality $g$ solves
\[
\max_{\{\bar{H}_i^R,\bar{H}_i^F\}_{i \in S^g}}\; W^g = \sum_{i \in S^g} N_i\,(\bar{U}-\phi(N_i, L_i))
\]
where $N_i$ is the equilibrium population at location $i$ and $\bar{U}_i = \Gamma\!\left(1-\frac{1}{\epsilon}\right)\Phi_i^{1/\epsilon}$ is the expected utility of a representative resident at location $i$, subject to $\bar{H}_i^R + \bar{H}_i^F = \bar{H}_i$ for each $i \in S^g$, taking other municipalities' allocations $\{\bar{H}_{i'}^R, \bar{H}_{i'}^F\}_{i' \notin S^g}$ and transportation investment $\{I_{jk}\}$ as given.

Substituting $\bar{H}_i^F = \bar{H}_i - \bar{H}_i^R$ and differentiating $\sum_{i \in S^g} N_i\bar{U}_i$ with respect to $\bar{H}_i^R$, the interior first-order condition at location $i \in S^g$ is
\begin{align}
\underbrace{\bar{U}_i\,\frac{\partial N_i}{\partial \bar{H}_i^R}}_{\geq\;0\;\text{(in-movers)}}
+\;
N_i\!\underbrace{\frac{\partial \bar{U}_i}{\partial q_i^R}\frac{\partial q_i^R}{\partial \bar{H}_i^R}}_{\geq\;0\;\text{(rent benefit)}}
+\;
N_i\!\underbrace{\frac{\partial \bar{U}_i}{\partial \phi}\frac{\partial \phi}{\partial N_i}\frac{\partial N_i}{\partial \bar{H}_i^R}}_{\leq\;0\;\text{(congestion cost)}}
+\;
N_i\!\underbrace{\left(-\frac{\partial \bar{U}_i}{\partial \phi}\frac{\partial \phi}{\partial L_i}\frac{\partial L_i}{\partial \bar{H}_i^F}\right)}_{\geq\;0\;\text{(less worker influx)}}
= 0.
\end{align}
The first term is positive: looser zoning lowers rents, attracting additional residents from the rest of the metro area; each in-mover contributes $\bar{U}_i$ to the municipality's objective. The second term is positive: lower rents raise effective income for all $N_i$ existing residents. The third term is negative: in-movers raise $N_i$, worsening congestion for all current residents through $\phi$. The fourth term is positive: shifting land from commercial to residential reduces employment $L_i$ and the associated worker congestion. The municipality sets $\bar{H}_i^R$ where the benefit of attracting residents and lowering rents is offset by the congestion they impose.

\subsection{Transportation Authority}

The metropolitan transportation authority chooses investment $\{I_{jk}\}_{j,k \in S}$ to maximize aggregate household welfare, taking the full zoning vector $\{\bar{H}_i^R, \bar{H}_i^F\}$ as given:
\[
\max_{\{I_{jk}\}}\; N \cdot \Gamma\!\left(1 - \tfrac{1}{\epsilon}\right)\Phi^{1/\epsilon}
\]
subject to the budget constraint
\[
\sum_{j,k} \delta_{jk}\, I_{jk} \leq K, \qquad I_{jk} \geq 0,
\]
where $\delta_{jk}$ is the per-unit cost of investment on link $(j,k)$ and $K$ is the total infrastructure budget.

Let $\mu \geq 0$ be the Lagrange multiplier on the budget constraint. The interior first-order condition for link $(j,k)$ is
\begin{align}
\frac{N}{\epsilon}\,\Gamma\!\left(1-\tfrac{1}{\epsilon}\right)\Phi^{1/\epsilon-1}
\frac{\partial\Phi}{\partial I_{jk}} = \mu\,\delta_{jk}.
\end{align}
The left-hand side is the marginal social benefit of investment on link $(j,k)$: lower commuting costs raise $\tilde{V}_{ij}$ for all commuters on that link, increasing $\Phi$ and hence aggregate welfare. The authority equates the marginal benefit per dollar of investment across all active links.

\section{Empirical Facts}

The model's first-order conditions generate three testable predictions. From the municipal zoning FOC (equation~142): locations with larger employment inflows face a larger marginal worker-influx cost from loosening residential zoning (\textbf{Proposition 1}), and locations with better transit access face additional upward pressure on population, which amplifies the congestion term and triggers a preemptive tightening of zoning (\textbf{Proposition 2}). From the transportation authority's FOC (equation~166): because the marginal social benefit of transit investment depends on ridership, which depends on the density of residents and jobs near the corridor, restrictive zoning that suppresses station-area density should lower the equilibrium return to transit investment (\textbf{Proposition 3}).

I test these three predictions using Dallas--Fort Worth (DFW) tract-level data and, for Proposition 3, a cross-metropolitan panel. These tests are best understood as \emph{equilibrium consistency checks} on the model's mechanism, not as causal estimates of the effect of zoning on transportation outcomes; a separate identification strategy, using historical transit routes, planned-route instruments, or policy shocks to zoning authority, would be needed to establish causality. The results are mixed: Propositions 1 and 2 are strongly supported within DFW, while Proposition 3 is not supported in its baseline cross-metro form and only weakly so in supplementary mechanism checks, which I report honestly below.

\subsection{Data and Zoning Proxy}

The primary proxy for zoning restrictiveness is the tract-level housing supply elasticity $\hat\gamma_i$ of \citet{Baum-Snow2024-ii}, available for the DFW metro area ($N=1{,}016$ tracts). Lower $\hat\gamma_i$ indicates that the local housing stock responds less to demand shocks, corresponding to a smaller $\bar H_i^R$ in the model. Employment inflows $L_i$ are constructed from LODES 2023 origin-destination flows, aggregated to the workplace tract. Transit access is measured using DART station locations (148 stations in DFW), with a tract classified as transit-adjacent if its centroid lies within 0.5 miles of a station. For the cross-metro test of Proposition 3, I use the Wharton Residential Land Use Regulatory Index (WRLURI) and existing fixed-guideway route-km from Transit Explorer, matched to Baum-Snow and Han metro-level elasticities through a manually checked metro-name crosswalk. As a robustness check on the zoning measure itself, I also use minimum-lot-area (MLA) estimates from \citet{Song2025-ft}'s codebook-based zoning estimator, which is a more direct, structurally grounded measure than either WRLURI or housing supply elasticity, though available only for a small, non-metro-representative sample of municipalities.

\subsection{Proposition 1: Employment Concentration Tightens Zoning}

The municipal FOC predicts $\partial \hat\gamma_i/\partial \log L_i < 0$: tracts with larger employment inflows impose a larger marginal worker-influx cost on incumbent residents, so municipalities respond with tighter residential zoning. I estimate
\[
\hat\gamma_i = \alpha + \beta_1 \log(1+L_i) + \beta_2\,\text{pctdis}_i + \varepsilon_i
\]
on the DFW tract sample ($N=1{,}016$), where $\text{pctdis}_i$ is the fractional distance from the tract centroid to the CBD, controlling for the urban density gradient.

\begin{table}[h]
\centering
\caption{Employment Concentration and Housing Supply Elasticity}
\label{tab:p1}
\begin{tabular}{lcccc}
\toprule
& (1) & (2) & (3) & (4) \\
& $\hat\gamma_{01a}$ & $\hat\gamma_{01a}$ & $\hat\gamma_{11a}$ & $\hat\gamma_{01a,\text{sp}}$ \\
\midrule
$\log(1+L_i)$ & $-0.009^{***}$ & $-0.009^{***}$ & $-0.006^{***}$ & $-0.009^{***}$ \\
 & (0.002) & (0.002) & (0.002) & (0.002) \\
$\text{pctdis}_i$ & & $-0.087^{*}$ & $-0.095^{**}$ & $-0.040$ \\
 & & (0.045) & (0.045) & (0.042) \\
\midrule
$R^2$ & 0.016 & 0.020 & 0.011 & 0.017 \\
$N$ & 1,016 & 1,016 & 1,016 & 1,016 \\
\bottomrule
\end{tabular}
\begin{minipage}{0.9\textwidth}
\footnotesize Note: Standard errors in parentheses. Columns (1)--(2) use the primary elasticity measure $\hat\gamma_{01a}$; column (3) uses the alternative 2010--2020 measure $\hat\gamma_{11a}$; column (4) uses a spatially smoothed version of $\hat\gamma_{01a}$. $^{*}p<0.10$, $^{**}p<0.05$, $^{***}p<0.01$.
\end{minipage}
\end{table}

The coefficient on $\log(1+L_i)$ is negative and significant at the 1\% level across all four specifications (Table~\ref{tab:p1}), and is robust to controlling for distance to the CBD and to using an alternative elasticity vintage. This is consistent with the worker-influx term in the municipal FOC: tracts that already attract substantial employment face tighter residential zoning, as incumbent residents weigh the congestion cost of additional in-commuters against the benefits of looser land use.

\subsection{Proposition 2: Transit Access Triggers a Zoning-Tightening Response}

The municipal FOC also implies that transit access, by raising residential attractiveness $\tilde V_{ij}$ and the resulting equilibrium population pressure $N_i$, should trigger a preemptive tightening of zoning: $\hat\gamma_i|_{\text{near transit}} < \hat\gamma_i|_{\text{far}}$. I estimate
\[
\hat\gamma_i = \alpha + \beta_1 \mathbf{1}[\text{near DART}]_i + \beta_2\,\text{pctdis}_i + \beta_3 \log(1+L_i) + \varepsilon_i
\]
on the subsample of $N=581$ Baum-Snow tracts matched to 2020 Census geometry, of which 15 have a DART-adjacent centroid.

\begin{table}[h]
\centering
\caption{DART Access and Housing Supply Elasticity}
\label{tab:p2}
\begin{tabular}{lcccc}
\toprule
& (1) & (2) & (3) & (4) \\
\midrule
$\mathbf{1}[\text{near DART}]$ & $-0.152^{**}$ & $-0.205^{***}$ & $-0.178^{***}$ & $-0.177^{**}$ \\
 & (0.059) & (0.058) & (0.057) & (0.079) \\
$\text{pctdis}$ & & $-0.412^{***}$ & $-0.384^{***}$ & $-0.384^{***}$ \\
 & & (0.059) & (0.059) & (0.059) \\
$\log(1+L_i)$ & & & $-0.032^{***}$ & $-0.032^{***}$ \\
 & & & (0.007) & (0.007) \\
$\mathbf{1}[\text{near}]\times\text{pctdis}$ & & & & $-0.011$ \\
 & & & & (0.418) \\
\midrule
$R^2$ & 0.011 & 0.087 & 0.118 & 0.118 \\
$N$ & 581 & 581 & 581 & 581 \\
\bottomrule
\end{tabular}
\begin{minipage}{0.9\textwidth}
\footnotesize Note: Standard errors in parentheses. Sample restricted to the 581 Baum-Snow tracts with matched 2020 Census geometry. $^{*}p<0.10$, $^{**}p<0.05$, $^{***}p<0.01$.
\end{minipage}
\end{table}

DART-adjacent tracts have housing supply elasticities roughly 0.15--0.21 units lower than non-adjacent tracts, significant at the 1--5\% level across specifications (Table~\ref{tab:p2}), and the result is robust to controlling separately for CBD distance and employment concentration. The interaction between transit proximity and CBD distance is not significant, likely reflecting the small number of suburban DART-adjacent tracts in the sample. Taken together with Proposition 1, this provides within-DFW support for the model's congestion channel: both employment concentration and transit access, the two forces that raise local population pressure in the model, predict tighter observed zoning.

\subsection{Proposition 3: Restrictive Zoning and the Return to Transit Investment}

Proposition 3 is the least well supported of the three predictions, and I report this honestly. The transportation authority's FOC implies that the marginal social benefit of investment on a link rises with ridership, which in turn depends on residential and employment density near the corridor. If zoning suppresses that density, the model predicts a lower equilibrium return to transit investment, and hence, across metropolitan areas, a smaller transit network relative to population in more restrictively zoned metros.

\subsubsection{Cross-Metro Test}

The baseline specification regresses log existing fixed-guideway transit route-km on WRLURI and controls, across a panel of 41 U.S. metros matched between Transit Explorer and Baum-Snow and Han:
\[
\log(\text{TransitKm}_m) = \alpha + \beta_1 \text{WRLURI}_m + \beta_2 \log(\text{numtracts}_m) + \beta_3 \text{FracUnavail}_m + \varepsilon_m.
\]
The model predicts $\hat\beta_1 < 0$.

\begin{table}[h]
\centering
\caption{WRLURI and Transit Route-Km Across Metros}
\label{tab:p3-crossmetro}
\begin{tabular}{lccccc}
\toprule
& (1) & (2) & (3) & (4) & (5) \\
& Full & Postwar & +Year & B\&H & B\&H, postwar \\
\midrule
WRLURI & $0.687^{*}$ & $0.767^{*}$ & $0.732^{*}$ & & \\
 & (0.351) & (0.440) & (0.371) & & \\
$\text{elast\_inner}$ & & & & $-1.888$ & $-1.613$ \\
 & & & & (1.992) & (2.502) \\
$\log(\text{numtracts})$ & $1.001^{***}$ & $0.949^{**}$ & & $1.001^{***}$ & \\
 & (0.264) & (0.366) & & (0.264) & \\
\midrule
$R^2$ & 0.427 & 0.307 & 0.430 & 0.383 & 0.245 \\
$N$ & 41 & 33 & 41 & 41 & 33 \\
\bottomrule
\end{tabular}
\begin{minipage}{0.95\textwidth}
\footnotesize Note: Standard errors in parentheses. ``Postwar'' restricts to metros whose first fixed-guideway transit opened after 1945, dropping eight legacy transit metros (Boston, Philadelphia, New York, Pittsburgh, Chicago, Washington, Cleveland, San Francisco Bay Area). Column (3) adds the opening year of the first line as a control. Columns (4)--(5) replace WRLURI with the Baum-Snow and Han metro-level inner-ring elasticity. $^{*}p<0.10$, $^{**}p<0.05$, $^{***}p<0.01$.
\end{minipage}
\end{table}

The result does not support the model's prediction: WRLURI enters with the \emph{wrong sign}, positive rather than negative, and is significant at the 10\% level in all three specifications, including after controlling for legacy pre-1945 transit status and restricting to the postwar subsample (Table~\ref{tab:p3-crossmetro}, columns 1--3). Replacing WRLURI with the Baum-Snow and Han metro-level elasticity (columns 4--5) recovers the theoretically correct negative sign, but the coefficient is never statistically distinguishable from zero. The most plausible interpretation is that metro-level zoning restrictiveness is confounded with economic dynamism and historical transit investment: metros that are old, dense, and transit-rich often also developed stringent land-use regulation over the twentieth century, so a static cross-section cannot separate the model's mechanism from this reverse historical pattern.

\subsubsection{Mechanism Check: Station-Area Density}

Because the cross-metro test conflates current zoning with transit history, I test the model's underlying density channel more directly using within-DFW, station-area variation. I regress tract-level population and housing-unit density on the Baum-Snow and Han elasticity, DART proximity, and their interaction, for the $N=590$ matched DFW tracts:
\[
\log(\text{density}_i) = \alpha + \beta_1 \hat\gamma_i + \beta_2\,\mathbf{1}[\text{near DART}]_i + \beta_3\, \hat\gamma_i \times \mathbf{1}[\text{near DART}]_i + X_i'\theta + \varepsilon_i.
\]
If permissive zoning supports higher station-area density, the model implies $\beta_1 > 0$, or at least $\beta_1+\beta_3>0$ near transit.

\begin{table}[h]
\centering
\caption{Station-Area Density and Housing Supply Elasticity (Dallas)}
\label{tab:p3-density}
\begin{tabular}{lcc}
\toprule
& Log pop. density & Log unit density \\
\midrule
$\hat\gamma_i$ & $-2.295^{***}$ & $-2.549^{***}$ \\
 & (0.133) & (0.132) \\
$\mathbf{1}[\text{near DART}]$ & $-0.923^{***}$ & $-0.841^{***}$ \\
 & (0.231) & (0.229) \\
$\hat\gamma_i \times \mathbf{1}[\text{near DART}]$ & $1.747^{**}$ & $1.507^{*}$ \\
 & (0.786) & (0.781) \\
Controls & Yes & Yes \\
\midrule
$R^2$ & 0.488 & 0.512 \\
$N$ & 590 & 590 \\
\bottomrule
\end{tabular}
\begin{minipage}{0.85\textwidth}
\footnotesize Note: Standard errors in parentheses. ``Near DART'' indicates a tract centroid within 0.5 miles of a DART station; controls include CBD distance. $^{*}p<0.10$, $^{**}p<0.05$, $^{***}p<0.01$.
\end{minipage}
\end{table}

This check also fails to confirm the model's channel, though in a more informative way. The direct relationship between elasticity and density is negative, not positive (Table~\ref{tab:p3-density}): in DFW, tracts with \emph{lower} (more restrictive) elasticity are denser, most plausibly because central, high-demand, transit-accessible locations are simultaneously denser and more politically constrained, an equilibrium pattern the static cross-section cannot separate from the model's causal channel. The interaction with DART proximity is positive and significant, indicating that the negative elasticity-density gradient is attenuated near transit, but the implied near-DART slope remains negative ($-2.295+1.747=-0.548$ for population density). The evidence is therefore consistent with substantial confounding rather than with the model's density channel.

\subsubsection{Robustness: Direct Zoning Measures from Song (2025)}

Both WRLURI and the Baum-Snow and Han elasticity are indirect proxies for zoning restrictiveness, estimated from observed housing supply responses rather than from the underlying land-use rules. As a robustness check, I use minimum-lot-area (MLA) estimates from \citet{Song2025-ft}'s codebook-based zoning estimator, which measures actual minimum lot size by zoning district. Coverage is limited to a small, non-metro-representative sample: \citet{Song2025-ft}'s validation set, covering all Boston-area (MAPC) towns plus one sampled municipality per county in roughly 18 other counties, aggregated to 19 metros.

Re-estimating the cross-metro specification on the 10 metros common to both samples, only the direct Song MLA measure recovers the theoretically correct sign ($\hat\beta = -0.39$ to $-0.76$, bivariate and with controls), while WRLURI remains wrong-signed even in this small subsample and the Baum-Snow and Han elasticity is wrong-signed as well; none of the three is statistically significant at $N=10$. A separate validation regression of $\log(\text{MLA}_m)$ on WRLURI across 18 metros finds essentially no relationship ($\hat\beta = 0.005$, $p=0.97$), while the Baum-Snow and Han elasticity has the theoretically expected sign but is statistically indistinguishable from zero ($\hat\beta=-0.40$, $p=0.59$). A metro-level station-area density mechanism check using the same MLA data ($N=11$ metros) finds coefficients on $\log(\text{MLA})$ that are small, unstable in sign, and far from significant.

These supplementary checks do not confirm Proposition 3, but they are informative about \emph{why} the baseline cross-metro test fails: WRLURI shows essentially zero correlation with a direct, codebook-based measure of zoning stringency for the same places, suggesting that the wrong-signed WRLURI result partly reflects a measurement-validity problem with WRLURI itself, rather than a failure of the zoning-transit complementarity that the model predicts. All of the Song-based checks, however, are limited to $N=10$--19 metros drawn from a non-representative validation sample, and should be read as suggestive rather than conclusive.

\subsection{Summary}

Propositions 1 and 2, both implications of the municipal zoning FOC tested within DFW, are strongly supported: employment-concentrated tracts and DART-adjacent tracts are both significantly more restrictively zoned, consistent with the model's congestion and transit-pressure channels. Proposition 3, the transportation authority's FOC implication that restrictive zoning lowers the return to transit investment, is not supported by the baseline cross-metro WRLURI test, which is wrong-signed; the DFW station-area density mechanism check likewise does not confirm the channel, though it is consistent with substantial confounding between zoning, density, and transit history rather than a clean rejection. Supplementary checks using a direct, codebook-based zoning measure suggest that part of the difficulty lies in WRLURI's validity as a proxy for the specific land-use lever, minimum lot size or density caps, through which the model's mechanism operates. Taken together, these results provide reasonable empirical discipline for the municipal side of the model, while indicating that a credible test of the transportation authority's response to zoning will require either a more direct, plausibly exogenous zoning measure or a panel/event-study design exploiting the timing of transit investment, rather than a static cross-metro comparison.

\printbibliography

\end{document}